\documentclass[11pt,a4paper]{article}
\usepackage[utf8]{inputenc}
\usepackage[french]{babel}
\usepackage[T1]{fontenc}
\usepackage{lmodern}
\usepackage{amsmath}
\usepackage{amssymb}
\usepackage{amsfonts}
\usepackage{graphicx}
\usepackage{geometry}
\usepackage{booktabs}
\usepackage{float}
\usepackage{hyperref}
\usepackage{xcolor}
\usepackage{microtype}
\usepackage{pgfplots}
\pgfplotsset{compat=1.18}

\geometry{top=2.5cm, bottom=2.5cm, left=2.5cm, right=2.5cm}

\definecolor{codeblue}{rgb}{0.1,0.2,0.6}
\definecolor{darkgreen}{rgb}{0.1,0.5,0.1}
\definecolor{barblue}{rgb}{0.2,0.4,0.8}
\definecolor{linegreen}{rgb}{0.1,0.6,0.3}

\hypersetup{
    colorlinks=true,
    linkcolor=codeblue,
    filecolor=magenta,
    urlcolor=codeblue,
    pdftitle={Régulation de feux par DQN sous SUMO - Évaluation contre baselines},
    pdfpagemode=FullScreen,
}

\title{\textbf{Régulation adaptative d'un carrefour par Deep Q-Learning} \\
\Large Entraînement, évaluation contre baselines et analyse critique \\
\normalsize Carrefour réel de Thiès (Sénégal) sous Eclipse SUMO}

\author{\textbf{Mamadou Ahmadou Ka} \\
Étudiant en Licence 1 Informatique \\
Département Informatique \\
Université Iba Der Thiam de Thiès (UIDT), Sénégal}

\date{8 Août 2026}

\begin{document}

\maketitle

\begin{abstract}
Ce document présente l'entraînement et l'évaluation d'un agent de Deep Q-Learning
régulant un carrefour réel de Thiès (Sénégal) sous le simulateur microscopique
Eclipse SUMO. L'agent a été entraîné sur 150 épisodes ($108\,000$ décisions), la
demande de trafic étant régénérée à chaque épisode avec une graine différente afin
de contraindre la généralisation.

Le résultat central n'est pas la progression de l'agent par rapport à son
initialisation aléatoire --- qui ne prouve rien --- mais sa comparaison à deux
contrôleurs de référence sur cinq graines de trafic jamais vues à l'entraînement.
Face au feu à programme fixe actuellement en place, le régulateur appris réduit le
temps perdu par véhicule de $50{,}7\,\%$ ($31{,}5\text{ s} \rightarrow 15{,}5\text{ s}$)
et le temps d'attente de $62{,}2\,\%$ ($23{,}8\text{ s} \rightarrow 9{,}0\text{ s}$).
Un résultat contre-intuitif est également rapporté et expliqué : la baseline
actionnée par détection de véhicules se révèle \emph{moins} performante que le feu
fixe ($+8{,}7\,\%$ de temps perdu) sur ce carrefour à demande asymétrique.

Le rapport documente enfin les limites méthodologiques de l'expérimentation, dont
un défaut de sélection du modèle détecté et corrigé en cours d'étude.
\end{abstract}

\newpage
\tableofcontents
\newpage

\section{Introduction et Contexte Expérimental}

La régulation de trafic par apprentissage par renforcement repose sur l'interaction
entre un agent décisionnel et un environnement de simulation microscopique. L'étude
porte sur le carrefour
\texttt{cluster\_6021368912\_6083287698\_6083287699\_6083287700} de la ville de
Thiès, comportant quatre phases vertes et huit voies entrantes contrôlées.

\subsection{Formulation du problème de décision markovien}

\paragraph{État.} À chaque décision, l'agent observe un vecteur
$\mathbf{s}_t \in \mathbb{R}^{13}$ composé de :
\begin{itemize}
  \item 4 longueurs de file d'attente, agrégées par branche et normalisées ;
  \item 4 densités d'occupation, rapportées à la capacité de chaque groupe de voies ;
  \item l'encodage \textit{one-hot} de la phase verte courante (4 composantes) ;
  \item 1 indicateur binaire signalant que la durée de vert minimale est écoulée.
\end{itemize}

\paragraph{Action.} L'action $a_t \in \{0,1,2,3\}$ désigne \textbf{l'indice de la
prochaine phase verte}. Ce choix est important : un espace binaire
$\{\text{STAY}, \text{SWITCH}\}$ ne permettrait à l'agent que de suivre ou d'accélérer
un cycle prédéfini, sans jamais réordonner les phases. Les actions incompatibles avec
les contraintes de sécurité (vert minimal de $10\text{ s}$, vert maximal de $60\text{ s}$)
sont masquées : leur valeur $Q$ est forcée à $-\infty$ avant la sélection.

\paragraph{Récompense.} $r_t = (W_{t-1} - W_t)/100$, où $W$ désigne le temps
d'attente cumulé de tous les véhicules sur les voies contrôlées. La récompense est
positive lorsque l'agent résorbe l'attente. Le coût d'un changement de phase est
capturé naturellement par le temps de jaune, sans pénalité arbitraire.

\paragraph{Invariant temporel.} Le jaune ($3\text{ s}$) est \emph{inclus} dans
l'intervalle de décision ($5\text{ s}$). Rester sur la phase courante et en changer
consomment donc exactement la même durée simulée, ce qui rend les récompenses des
deux branches comparables. Sans cet invariant, changer de phase coûterait
mécaniquement plus de temps simulé et l'agent apprendrait à ne jamais commuter --- un
biais qui a effectivement été observé dans une version antérieure de
l'implémentation.

\subsection{Contrôle effectif du feu}

L'agent commande le carrefour par \texttt{setRedYellowGreenState()}, et non par
\texttt{setPhase()}. Cette distinction n'est pas cosmétique : \texttt{setPhase()}
positionne la phase mais réarme sa durée programmée, si bien que SUMO poursuit
ensuite le cycle statique de lui-même. L'agent croit alors piloter un feu qui, en
réalité, continue de tourner seul.

Un test dédié (\texttt{rl.smoke\_test}) vérifie cette propriété : le feu est laissé
$60\text{ s}$ sans aucun ordre et doit conserver son état, là où le programme statique
aurait basculé après $27\text{ s}$. Un garde-fou (\texttt{assert\_in\_sync}) compare
en outre l'état réel du feu à l'état interne de l'agent.

\section{Protocole expérimental}

\begin{table}[H]
\centering
\caption{Paramètres de l'expérimentation}
\label{tab:protocole}
\vspace{0.2cm}
\begin{tabular}{ll}
\toprule
\textbf{Élément} & \textbf{Valeur} \\
\midrule
Algorithme & Double DQN, perte de Huber, cible en \textit{soft update} ($\tau = 0{,}005$) \\
Architecture & Perceptron multicouche $13 \rightarrow 64 \rightarrow 64 \rightarrow 4$ \\
Optimiseur & Adam, $\alpha = 5 \times 10^{-4}$, $\gamma = 0{,}99$ \\
Tampon d'expérience & $50\,000$ transitions, lots de $64$ \\
Début d'apprentissage & après $1\,000$ transitions collectées \\
Exploration & $\epsilon$ linéaire de $1{,}0$ à $0{,}05$ sur $30\,000$ pas \\
\midrule
Durée d'un épisode & $3\,600\text{ s}$ simulées $\rightarrow$ 720 décisions \\
Nombre d'épisodes & 150 ($108\,000$ décisions) \\
Graines d'entraînement & $1000 + n$ pour l'épisode $n$ (demande régénérée) \\
Graines d'évaluation & $\{7, 13, 21, 42, 77\}$, \textbf{jamais vues} \\
\bottomrule
\end{tabular}
\end{table}

\paragraph{Demande de trafic.} La demande OSM d'origine laissait le carrefour en
écoulement libre ($469\text{ véh/h}$ dispersés sur toute la ville, dont seulement 12
trajets traversant les branches étudiées) : il n'y avait rien à optimiser. Une
demande dédiée a donc été construite, concentrée sur les quatre branches, avec un
axe majeur chargé deux fois plus que l'axe secondaire et un profil temporel en trois
paliers (montée, pointe, décrue). Cette asymétrie est ce qui rend le problème de
régulation non trivial.

\paragraph{Régénération à chaque épisode.} La demande est régénérée avec une graine
différente à chaque épisode. L'agent ne peut donc pas mémoriser une séquence
d'arrivées particulière ; il doit apprendre une politique fonction de l'état observé.

\section{Résultats d'entraînement}

Le Tableau~\ref{tab:metrics_summary} présente l'évolution des indicateurs internes
au fil de l'entraînement.

\begin{table}[H]
\centering
\caption{Évolution des indicateurs d'entraînement (moyennes par tranche de 25 épisodes)}
\label{tab:metrics_summary}
\vspace{0.2cm}
\begin{tabular}{lcccccc}
\toprule
\textbf{Épisodes} & \textbf{1--25} & \textbf{26--50} & \textbf{51--75} & \textbf{76--100} & \textbf{101--125} & \textbf{126--150} \\
\midrule
Attente moyenne ($s$)          & $413{,}3$ & $91{,}0$ & $50{,}5$ & $43{,}2$ & $38{,}9$ & $\mathbf{34{,}1}$ \\
File moyenne (\textit{véh})    & $8{,}35$  & $4{,}31$ & $3{,}13$ & $2{,}85$ & $2{,}65$ & $\mathbf{2{,}48}$ \\
Vitesse sur voies ($km/h$)     & $33{,}6$  & $40{,}7$ & $45{,}5$ & $47{,}6$ & $48{,}6$ & $\mathbf{49{,}4}$ \\
\bottomrule
\end{tabular}
\end{table}

\begin{table}[H]
\centering
\caption{Valeurs aux jalons de l'entraînement}
\label{tab:jalons}
\vspace{0.2cm}
\begin{tabular}{lcccc}
\toprule
\textbf{Métrique} & \textbf{Épisode 1} & \textbf{Épisode 50} & \textbf{Épisode 100} & \textbf{Épisode 150} \\
\midrule
Attente moyenne ($s$)         & $326{,}4$ & $51{,}5$   & $41{,}5$   & $\mathbf{27{,}0}$ \\
File moyenne (\textit{véh})   & $8{,}45$  & $3{,}34$   & $2{,}70$   & $\mathbf{2{,}16}$ \\
Perte de Huber                & $0{,}0$   & $0{,}0859$ & $0{,}0158$ & $\mathbf{0{,}0077}$ \\
Exploration $\epsilon$        & $0{,}9772$& $0{,}0500$ & $0{,}0500$ & $0{,}0500$ \\
Changements de phase          & $213$     & $216$      & $221$      & $218$ \\
\bottomrule
\end{tabular}
\end{table}

\paragraph{Avertissement sur la métrique de vitesse.} La « vitesse sur voies » est la
moyenne de \texttt{getLastStepMeanSpeed} sur les voies contrôlées. Sur une voie
\emph{vide}, SUMO renvoie la vitesse autorisée. Cet indicateur mélange donc la vitesse
réelle du trafic et la proportion de voies dégagées, ce qui explique qu'il approche
$51\text{ km/h}$ alors que la limite urbaine est de $50\text{ km/h}$. Il est
utilisable comme signal de progression relative, mais \textbf{ne doit pas être
présenté comme la vitesse des véhicules}. La vitesse réelle, calculée par véhicule à
partir de \texttt{tripinfo.xml}, figure au Tableau~\ref{tab:baselines} et vaut
$8{,}16\text{ m/s}$, soit $29{,}4\text{ km/h}$.

\subsection{Progression de l'attente moyenne}

\begin{figure}[H]
\centering
\begin{tikzpicture}
\begin{axis}[
    width=0.88\textwidth,
    height=6.5cm,
    xlabel={Épisodes d'entraînement},
    ylabel={Attente moyenne ($s$)},
    xmin=1, xmax=150,
    ymin=0, ymax=400,
    grid=major,
    grid style={dashed,gray!30},
    thick,
    axis lines=left,
    legend pos=north east
]
\addplot[color=codeblue, mark=*, mark size=1.5pt, thick] coordinates {
    (1,326.4) (10,348.7) (20,156.0) (30,132.8) (40,71.7) (50,51.5)
    (60,46.1) (70,50.1) (80,36.1) (90,35.6) (100,41.5) (110,36.7)
    (120,36.4) (130,39.8) (140,30.1) (150,27.0)
};
\addlegendentry{Attente moyenne ($s$)}
\end{axis}
\end{tikzpicture}
\caption{Effondrement de l'attente moyenne, avec plateau atteint vers l'épisode 70.}
\label{fig:wait_curve}
\end{figure}

\subsection{Décroissance de la perte}

La perte optimisée est la \textbf{perte de Huber} (\texttt{smooth\_l1\_loss}), moins
sensible aux valeurs aberrantes que l'erreur quadratique --- propriété utile ici, les
récompenses étant occasionnellement extrêmes lors des blocages.

\begin{figure}[H]
\centering
\begin{tikzpicture}
\begin{axis}[
    width=0.88\textwidth,
    height=6.5cm,
    xlabel={Épisodes d'entraînement},
    ylabel={Perte de Huber},
    xmin=1, xmax=150,
    ymin=0, ymax=0.55,
    grid=major,
    grid style={dashed,gray!30},
    thick,
    axis lines=left,
    legend pos=north east
]
\addplot[color=linegreen, mark=square*, mark size=1.5pt, thick] coordinates {
    (2,0.40751) (10,0.38072) (20,0.36461) (30,0.29297) (40,0.07248)
    (50,0.08586) (60,0.10381) (70,0.10461) (80,0.02757) (90,0.0396)
    (100,0.01582) (110,0.01739) (120,0.0145) (130,0.01457)
    (140,0.01588) (150,0.00772)
};
\addlegendentry{Perte de Huber}
\end{axis}
\end{tikzpicture}
\caption{Décroissance de la perte du réseau. La valeur nulle à l'épisode 1 traduit
l'absence de mise à jour tant que le tampon n'a pas atteint $1\,000$ transitions.}
\label{fig:loss_curve}
\end{figure}

\subsection{Distribution des attentes moyennes}

\begin{figure}[H]
\centering
\begin{tikzpicture}
\begin{axis}[
    ybar,
    bar width=18pt,
    width=0.88\textwidth,
    height=6.5cm,
    ylabel={Nombre d'épisodes},
    symbolic x coords={[<=50s], [51-100s], [101-300s], [301-600s], [>600s]},
    xtick=data,
    nodes near coords,
    nodes near coords align={vertical},
    ymin=0, ymax=90,
    grid=major,
    grid style={dashed,gray!30},
    axis lines=left,
    enlarge x limits=0.15
]
\addplot[fill=barblue!80, draw=codeblue] coordinates {
    ([<=50s],78)
    ([51-100s],38)
    ([101-300s],20)
    ([301-600s],10)
    ([>600s],4)
};
\end{axis}
\end{tikzpicture}
\caption{Répartition des attentes moyennes sur les 150 épisodes. Les 14 épisodes
au-delà de $300\text{ s}$ appartiennent tous à la phase d'exploration initiale.}
\label{fig:wait_histogram}
\end{figure}

\section{Évaluation contre baselines}

\subsection{Pourquoi la progression interne ne prouve rien}

Comparer l'agent entraîné à l'agent de l'épisode 1 revient à le comparer à un réseau
de neurones \emph{aléatoire}. Tout agent qui apprend quoi que ce soit bat son
initialisation ; un tel écart ne renseigne donc pas sur l'utilité pratique du
système. La seule question qui engage le projet est : \textbf{le régulateur appris
fait-il mieux que le feu réellement en place ?}

Deux contrôleurs de référence sont donc évalués :

\begin{itemize}
  \item \textbf{fixe} --- le programme statique d'origine du réseau
        ($27/6/6/6/27/6/6/6$), c'est-à-dire le feu tel qu'il fonctionne aujourd'hui ;
  \item \textbf{actué} --- un feu à détection de véhicules, obtenu par
        \texttt{netconvert --tls.default-type actuated --tls.rebuild}
        ($minDur = 8\text{ s}$, $maxDur = 50\text{ s}$). Le vert est prolongé tant que
        des véhicules se présentent. C'est la baseline exigeante.
\end{itemize}

Le type actionné du programme est vérifié à l'exécution : si le feu ne comporte aucun
programme de type \texttt{ACTUATED}, l'évaluation s'interrompt par une erreur. Sans ce
garde-fou, une baseline mal générée serait silencieusement identique au contrôleur
fixe, et le résultat publié serait faux.

\subsection{Résultats}

Les trois contrôleurs sont évalués sur les cinq graines réservées, avec la même
demande, sur des épisodes de $3\,600\text{ s}$. La politique apprise est exploitée
sans exploration ($\epsilon = 0$). Les métriques proviennent de
\texttt{tripinfo.xml}, donc mesurées \emph{par véhicule}.

\begin{table}[H]
\centering
\caption{Comparaison des trois contrôleurs (5 graines inédites, écart-type inter-graines)}
\label{tab:baselines}
\vspace{0.2cm}
\begin{tabular}{lccccc}
\toprule
\textbf{Contrôleur} & \textbf{Temps perdu ($s$)} & \textbf{Attente ($s$)} & \textbf{Trajet ($s$)} & \textbf{Vitesse ($m/s$)} & \textbf{$\Delta$} \\
\midrule
fixe   & $31{,}5 \pm 0{,}5$ & $23{,}8 \pm 0{,}4$ & $48{,}5$ & $6{,}37$ & référence \\
actué  & $34{,}2 \pm 1{,}1$ & $25{,}8 \pm 0{,}8$ & $51{,}2$ & $5{,}72$ & $+8{,}7\,\%$ \\
\textbf{DQN} & $\mathbf{15{,}5 \pm 0{,}6}$ & $\mathbf{9{,}0 \pm 0{,}5}$ & $\mathbf{32{,}7}$ & $\mathbf{8{,}16}$ & $\mathbf{-50{,}7\,\%}$ \\
\bottomrule
\end{tabular}
\end{table}

Le régulateur appris divise le temps perdu par deux et l'attente par $2{,}6$ face au
feu en place. Entre $843$ et $907$ véhicules terminent leur trajet selon la graine, avec
moins de $1\,\%$ d'écart entre contrôleurs : les trois écoulent bien la même demande,
et la comparaison porte donc sur la qualité de la régulation, non sur un débit
différent.

\subsection{Le gain moyen cache-t-il des véhicules sacrifiés ?}
\label{sec:queues}

Une moyenne peut s'améliorer en dégradant une minorité : un contrôleur qui fluidifie
l'axe majeur au prix de quelques véhicules bloqués très longtemps sur l'axe secondaire
afficherait un bon score moyen tout en étant inacceptable. La question se vérifie sur
la distribution complète, et non sur la moyenne seule ($4\,386$ véhicules pour l'agent,
$4\,360$ pour le contrôleur fixe, cumulés sur les cinq graines).

\begin{table}[H]
\centering
\caption{Distribution du temps d'attente par véhicule (quantiles, en secondes)}
\label{tab:quantiles}
\vspace{0.2cm}
\begin{tabular}{lcccccc}
\toprule
\textbf{Contrôleur} & \textbf{$n$} & \textbf{moyenne} & \textbf{médiane} & \textbf{$p_{90}$} & \textbf{$p_{99}$} & \textbf{maximum} \\
\midrule
fixe   & $4\,360$ & $23{,}8$ & $20$ & $55$ & $89$ & $127$ \\
actué  & $4\,375$ & $25{,}8$ & $24$ & $53$ & $76$ & $136$ \\
\textbf{DQN} & $4\,386$ & $\mathbf{9{,}0}$ & $\mathbf{5}$ & $\mathbf{25}$ & $\mathbf{45}$ & $\mathbf{89}$ \\
\bottomrule
\end{tabular}
\end{table}

L'agent domine à \emph{tous} les quantiles, y compris au maximum : le véhicule le plus
pénalisé sous sa régulation attend $89\text{ s}$, contre $127\text{ s}$ sous le feu
fixe et $136\text{ s}$ sous le feu actionné. Le gain n'est donc pas obtenu par
report sur une minorité --- le pire cas s'améliore de $30\,\%$ en même temps que la
moyenne. C'est une propriété qui ne se déduisait pas du Tableau~\ref{tab:baselines} et
qui devait être vérifiée séparément.

\subsection{Un résultat contre-intuitif : la baseline actionnée est moins bonne que le feu fixe}

Le feu actionné dégrade les performances de $8{,}7\,\%$ par rapport au feu fixe. Ce
résultat surprend, la détection de véhicules étant réputée supérieure à un programme
statique. Il s'explique par la structure de la demande.

La logique actionnée prolonge le vert d'une branche tant que des véhicules s'y
présentent, sur un critère purement \emph{local}. Or l'axe majeur du carrefour est
chargé deux fois plus que l'axe secondaire : le vert de l'axe majeur est donc
systématiquement étendu jusqu'à $maxDur$, et l'axe secondaire est privé de vert. Le
programme fixe, lui, garantit un tour de rôle : moins réactif, mais équitable.

Cette faiblesse est précisément ce que le DQN exploite. Sa décision porte sur les
files observées \emph{sur les quatre branches simultanément}, ce qui lui permet
d'arbitrer entre elles plutôt que de réagir à une seule.

Ce résultat est conservé tel quel plutôt que d'être écarté : une baseline qui échoue
pour une raison identifiée est une information, et elle délimite le domaine de
validité de la conclusion --- le gain observé tient à l'asymétrie de la demande.

\section{Analyse critique et limites}

\subsection{Un défaut de sélection du modèle, détecté et corrigé}

La sauvegarde du « meilleur » modèle était initialement pilotée par le retour cumulé
de l'épisode. Ce critère est vide de sens dans cette formulation : la récompense
étant une différence $(W_{t-1} - W_t)/100$, sa somme sur un épisode se télescope en
\[
\sum_{t=1}^{T} \frac{W_{t-1} - W_t}{100} = \frac{W_0 - W_T}{100},
\]
soit une quantité qui ne dépend que des états initial et final --- du bruit autour de
zéro, ce que confirme la colonne \texttt{reward} du journal.

Conséquence mesurée : le modèle ainsi retenu (aux alentours de l'épisode 50) atteint
$-32{,}3\,\%$ de temps perdu, contre $-50{,}7\,\%$ pour le modèle de l'épisode 150.
Le critère a été remplacé par l'attente moyenne de l'épisode, évaluée uniquement une
fois $\epsilon$ au plancher (sinon l'exploration résiduelle fausse la comparaison).

Ce défaut est rapporté parce qu'il illustre un risque général : une métrique de
sélection non vérifiée peut faire publier un modèle sensiblement inférieur à celui
dont on dispose, sans qu'aucun signal d'erreur n'apparaisse.

\subsection{L'agent a-t-il simplement ré-appris un cycle fixe ?}
\label{sec:adaptivite}

Le nombre de changements par épisode se concentre dans $[188, 228]$, de moyenne
$\mu = 215{,}6$ --- une régularité qui suggère un cycle. L'argument par l'absurde
(un cycle fixe ne peut pas diviser par deux le temps perdu d'un cycle fixe) est
recevable mais indirect. La question est donc tranchée par la mesure directe
(\texttt{rl/adaptivity.py}, cinq graines d'évaluation, $1\,083$ verts observés pour
l'agent contre $795$ pour le contrôleur fixe).

\paragraph{Un test invalide, à écarter explicitement.} L'idée naturelle est de
corréler la durée d'un vert à la file de la branche desservie. Appliquée au
contrôleur \emph{fixe}, cette corrélation vaut $r = +0{,}54$ : un feu strictement
statique paraît donc « réagir » au trafic. L'explication est un confondant de
topologie --- les phases de l'axe majeur ont à la fois plus de vert ($33\text{ s}$
contre $12\text{ s}$) et plus de file. La corrélation mesure la structure du cycle,
pas une réaction à l'état. \textbf{Elle ne peut servir de preuve dans aucun sens} et
est écartée.

\paragraph{Test valide : la variance intra-phase.} Un cycle fixe a, par construction,
un écart-type \emph{rigoureusement nul} sur la durée d'une même phase. C'est sa
signature, et elle est falsifiable.

\begin{table}[H]
\centering
\caption{Modulation de la durée du vert, à phase constante (5 graines)}
\label{tab:adaptivite}
\vspace{0.2cm}
\begin{tabular}{lccc}
\toprule
 & \textbf{fixe} & \textbf{actué} & \textbf{DQN} \\
\midrule
Écart-type intra-phase & $\mathbf{0{,}00}$ s & $1{,}38$ s & $\mathbf{3{,}53}$ s \\
Plage des durées observées & $12$--$33$ s & $14$--$44$ s & $\mathbf{15}$--$\mathbf{55}$ s \\
Choix de la phase la plus chargée & --- & --- & $\mathbf{79{,}2\,\%}$ \\
\bottomrule
\end{tabular}
\end{table}

Le contrôleur fixe donne $0{,}00\text{ s}$, comme attendu. L'agent module la durée
d'une même phase de $15$ à $55\text{ s}$ selon l'état observé : \textbf{l'hypothèse du
cycle fixe est réfutée par la mesure}, et non seulement par l'absurde. Par ailleurs,
lorsque plusieurs actions sont autorisées, l'agent sélectionne la phase dont la file
est la plus longue dans $79{,}2\,\%$ des cas ($1\,208$ décisions), contre $25\,\%$ pour
un choix uniforme sur quatre phases.

La stabilité du \emph{nombre} de changements s'explique enfin par les contraintes de
sécurité : avec un vert minimal de $10\text{ s}$ et maximal de $60\text{ s}$ sur
$3\,600\text{ s}$, ce nombre est encadré par construction dans $[60, 360]$. La marge de
manœuvre porte sur le choix de la phase et l'instant du basculement, non sur la
fréquence globale.

\paragraph{Effet de plancher.} La file à l'ouverture du vert vaut en moyenne
$0{,}31$ véhicule chez l'agent, contre $2{,}24$ chez le contrôleur fixe. L'agent
maintient le carrefour si dégagé qu'il ne reste presque plus de variance à expliquer :
c'est un succès du contrôleur, mais cela rend tout test fondé sur cette variable peu
informatif --- raison supplémentaire de s'appuyer sur la variance intra-phase.

\subsection{Une phase abandonnée : optimisation réelle, et son coût}
\label{sec:phase1}

La mesure a révélé un comportement non anticipé : \textbf{l'agent n'ouvre jamais la
phase 1} ($0$ vert sur $1\,083$ observés ; $410$, $451$ et $222$ pour les phases $0$,
$2$ et $3$). L'analyse de la topologie du feu montre qu'il ne s'agit pas d'une
défaillance.

\begin{table}[H]
\centering
\caption{Mouvements desservis par chaque phase verte}
\label{tab:phases}
\vspace{0.2cm}
\begin{tabular}{clp{7.4cm}}
\toprule
\textbf{Phase} & \textbf{Liens verts} & \textbf{Mouvements} \\
\midrule
$0$ & $10$ & axe majeur tout droit (\texttt{G}) \emph{et} tourne-à-gauche en permissif (\texttt{g}) \\
$\mathbf{1}$ & $\mathbf{4}$ & \textbf{tourne-à-gauche de l'axe majeur uniquement, en protégé (\texttt{G})} \\
$2$ & $6$ & axe secondaire tout droit \\
$3$ & $4$ & tourne-à-gauche de l'axe secondaire, en protégé \\
\bottomrule
\end{tabular}
\end{table}

La phase 1 est un \textbf{sous-ensemble strict de la phase 0} : tous les mouvements
qu'elle autorise le sont déjà en phase 0, en permissif au lieu de protégé. L'agent a
donc identifié qu'elle n'apporte aucune capacité supplémentaire et redistribue ses
$\approx 12\text{ s}$ de vert aux mouvements qui en tirent parti. C'est une
optimisation de fond, que ni le contrôleur fixe ni le contrôleur actionné ne peuvent
découvrir, l'un comme l'autre parcourant les phases dans l'ordre programmé.

\paragraph{Le coût, qui doit être énoncé.} Faire passer un tourne-à-gauche du protégé
au permissif transfère la résolution du conflit au conducteur, qui doit trouver un
créneau dans le flux opposé. En simulation, le coût reste marginal ($2$ freinages
d'urgence sur cinq épisodes). Sur le terrain, ce compromis relève d'une décision
d'ingénierie de sécurité et \textbf{non de l'agent} : la fonction de récompense ne
mesure que le temps d'attente, et rien n'y représente le risque d'accident. Un agent
optimise ce qu'on lui donne à optimiser --- ici, il a trouvé un gain que le cahier des
charges implicite n'interdisait pas.

La correction appropriée n'est pas de brider l'agent après coup mais d'inscrire la
contrainte dans le problème : garantir un vert protégé périodique par mouvement, ou
pénaliser explicitement les conflits permissifs dans la récompense.

\subsection{Métrique de vitesse trompeuse sur les voies}

Comme signalé en Section~3, la vitesse moyenne par voie intègre les voies vides, où
SUMO renvoie la vitesse autorisée. Un carrefour parfaitement dégagé maximise donc
cette métrique, qui mesure en partie l'absence de trafic. Les conclusions du présent
rapport reposent pour cette raison sur \texttt{tripinfo.xml}, agrégé par véhicule.

\subsection{Équité par mouvement non garantie}

La récompense n'agrège que le temps d'attente \emph{total} : rien n'y interdit de
privilégier durablement un mouvement au détriment d'un autre. Le
Tableau~\ref{tab:quantiles} montre qu'aucun véhicule n'est sacrifié en pratique --- le
pire cas s'améliore lui aussi --- mais c'est une propriété \emph{constatée}, non une
propriété \emph{garantie} par la formulation. L'abandon de la phase 1
(Section~\ref{sec:phase1}) en est l'illustration : l'agent a supprimé un vert protégé
parce qu'aucun terme de la récompense ne le protégeait.

Une formulation destinée au déploiement devrait inscrire cette contrainte dans le
problème plutôt que d'espérer qu'elle émerge : plafond d'attente par mouvement,
garantie de service périodique, ou pénalité sur la variance inter-branches.

\subsection{Portée limitée à un carrefour isolé}

L'agent régule une intersection unique. Rien ne garantit que la politique apprise
reste bénéfique en réseau : fluidifier un carrefour peut déplacer la congestion vers
le suivant. La validation d'un déploiement réel exigerait une évaluation sur un
corridor, avec une coordination inter-carrefours.

\subsection{Absence de répétitions indépendantes}

Les résultats proviennent d'\emph{une} exécution d'entraînement évaluée sur cinq
graines. L'écart-type reporté mesure la variabilité \emph{du trafic}, non celle de
\emph{l'apprentissage}. Trois à cinq entraînements avec des graines d'initialisation
distinctes seraient nécessaires pour établir la stabilité de la méthode elle-même.
L'ampleur de l'écart observé ($-50{,}7\,\%$ contre $\pm 0{,}6\text{ s}$ d'écart-type)
rend toutefois improbable qu'il s'agisse d'un artefact.

\subsection{Écart entre simulation et terrain}

SUMO modélise un trafic homogène de véhicules particuliers. Le trafic réel de Thiès
comporte deux-roues, charrettes et piétons traversant hors passage, dont les
comportements ne sont pas représentés. Les gains rapportés valent dans le modèle ; ils
constituent une borne optimiste pour le terrain.

\section{Conclusion et perspectives}

\subsection{Rappels théoriques}

\begin{align}
  \textbf{Équation de Bellman optimale} \quad & Q^{*}(s,a)=\mathbb{E}_{s'}\bigl[ r + \gamma \max_{a'} Q^{*}(s',a') \mid s,a\bigr] \\
  \textbf{Cible du Double DQN} \quad & y_t = r_t + \gamma\, Q_{\bar\theta}\bigl(s_{t+1},\ \arg\max_{a'} Q_{\theta}(s_{t+1},a')\bigr)
\end{align}

Le découplage entre le réseau qui \emph{choisit} l'action ($\theta$) et celui qui
l'\emph{évalue} ($\bar\theta$) atténue la surestimation systématique des valeurs $Q$
propre au DQN classique.

\subsection{Résultats établis}

\begin{enumerate}
    \item Le régulateur appris réduit le temps perdu par véhicule de
          $\mathbf{50{,}7\,\%}$ face au feu fixe en place
          ($31{,}5\text{ s} \rightarrow 15{,}5\text{ s}$), sur cinq graines de trafic
          jamais rencontrées à l'entraînement.
    \item L'attente moyenne par véhicule diminue de $\mathbf{62{,}2\,\%}$
          ($23{,}8\text{ s} \rightarrow 9{,}0\text{ s}$).
    \item Le gain n'est pas obtenu au détriment d'une minorité : l'agent domine à
          \emph{tous} les quantiles, le véhicule le plus pénalisé attendant
          $89\text{ s}$ contre $127\text{ s}$ sous le feu fixe.
    \item Le caractère adaptatif est établi par la mesure et non par l'argument : la
          durée d'une même phase varie de $15$ à $55\text{ s}$ (écart-type intra-phase
          de $3{,}53\text{ s}$, contre $0{,}00\text{ s}$ pour le feu fixe), et l'agent
          sert la branche la plus chargée dans $79{,}2\,\%$ des arbitrages.
    \item Le contrôle effectif du feu par l'agent est vérifié par un test dédié, et
          non supposé.
    \item La baseline actionnée s'avère inférieure au feu fixe sur cette demande
          asymétrique, ce qui délimite le domaine de validité du résultat.
    \item L'agent a découvert que la phase 1 est redondante en capacité et l'a
          supprimée --- gain réel, mais au prix du passage de tourne-à-gauche protégés
          en permissifs, arbitrage qui relève de l'ingénieur et non de l'agent.
\end{enumerate}

\subsection{Travaux futurs}

\begin{itemize}
  \item Répéter l'entraînement avec plusieurs graines d'initialisation pour établir la
        stabilité de la méthode, et non seulement celle d'un modèle.
  \item Réintroduire l'équité par mouvement dans la formulation (vert protégé
        périodique garanti, ou pénalité sur les conflits permissifs) et vérifier le
        coût en performance de cette contrainte.
  \item Étendre à un corridor de plusieurs carrefours coordonnés, pour vérifier que le
        gain local ne se traduit pas par un report de congestion.
  \item Éprouver la robustesse sur des scénarios dégradés : pic de demande, panne de
        détecteur, arrivée de véhicule prioritaire.
  \item Publier poids (\texttt{dqn\_best.pt}), journal d'entraînement et scripts
        d'évaluation afin de permettre la reproduction indépendante des résultats.
\end{itemize}

\subsection{Reproductibilité}

\begin{verbatim}
py -3.14 -m rl.smoke_test                          # vérifie le contrôle du feu
py -3.14 -m rl.train --episodes 150                # entraînement (~45 min)
py -3.14 -m rl.evaluate --model outputs/dqn_best.pt  # tableau comparatif
py -3.14 -m rl.adaptivity --seeds 7 13 21 42 77    # preuve d'adaptativité
py -3.14 -m rl.demo --controller dqn --seed 7      # démonstration visuelle
\end{verbatim}

\paragraph{Avertissement de reproductibilité.} Les fichiers \texttt{tripinfo\_*.xml}
du répertoire \texttt{outputs/} sont écrasés à chaque exécution, et
\texttt{rl.train} réécrit \texttt{training\_log.csv} depuis zéro. Un artefact
présent sur le disque ne correspond donc pas nécessairement au modèle courant : lors
de la rédaction, des \texttt{tripinfo} antérieurs au remplacement de
\texttt{dqn\_best.pt} donnaient $-32{,}3\,\%$ au lieu de $-50{,}7\,\%$. Toute
vérification doit relancer \texttt{rl.evaluate} plutôt que relire les fichiers
existants.

\end{document}
